\section{Concentration energy and global dynamics at critical mass}
\label{sec:r11-variational}

A high-frequency concentration has a definite limiting energy contribution
from a bounded correction to the homogeneous dispersion.  We identify
that contribution using the full optimizer manifold.  Conservation then
converts the resulting compactness statement into a global theorem for
all initial data below the concentration energy, including data at the
critical mass.  No proximity to a standing wave is needed.

\subsection{The homogeneous problem and the concentration energy}

Fix $N\ge1$, $m>1$, $s=m/2$, $r=m/N$, and $p=2+2r$.
Let $a$ be real, even, positive elliptic, and homogeneous of degree $m$.
Write
\[
 M(f)=\|f\|_2^2,\qquad
 A(f)=\int a(\xi)|\widehat f(\xi)|^2\,d\xi,\qquad
 E_a(f)=\tfrac12 A(f)-\tfrac1p\|f\|_p^p.
\]
Let $M_c$ be the sharp mass in
\begin{equation}\label{eq:r11-sharp-gn}
 \frac1p\|f\|_p^p\le
 \frac12\left(\frac{M(f)}{M_c}\right)^r A(f),
\end{equation}
and define the normalized equality manifold
\[
 \mathcal G=\{Q\in H^s:M(Q)=M_c,\ A(Q)=1,\ E_a(Q)=0\}.
\]
We use the normalized sharp-GN compactness property: whenever
$M(f_n)\to M_c$, $A(f_n)\to1$, and $E_a(f_n)\to0$, there are
translations and a subsequence converging strongly in $H^s$ to an
element of $\mathcal G$.  For the homogeneous symbols considered in
this work this is the established sharp-GN compactness input.  It
requires neither optimizer uniqueness nor nondegeneracy of the real
linearized operator.  In particular $\mathcal G$ is compact modulo
translations; constant phases are already included in this set.

Let $d$ be a bounded, measurable, real, even Fourier multiplier. Assume
that there is a bounded function $d_\infty$ on $S^{N-1}$ such that
\begin{equation}\label{eq:r11-angular-limit}
 \sup_{\theta\in S^{N-1}}|d(R\theta)-d_\infty(\theta)|
       \longrightarrow0\qquad(R\longrightarrow\infty).
\end{equation}
Define
\[
 k=a-d,\qquad
 D(f)=\int d(\xi)|\widehat f(\xi)|^2\,d\xi,\qquad
 E_k(f)=E_a(f)-\tfrac12D(f).
\]
No sign assumption is imposed on $d$ or $d_\infty$.  The quadratic
perturbation is continuous on $L^2$, and $E_k$ is continuous on $H^s$.
Set
\begin{equation}\label{eq:r11-concentration-energy}
 e_\infty=-\frac12\max_{Q\in\mathcal G}
       \int d_\infty(\xi/|\xi|)|\widehat Q(\xi)|^2\,d\xi.
\end{equation}
The maximum exists: its objective is a bounded-multiplier quadratic
form, continuous in $L^2$ and invariant under translations and constant
phase.  Compactness of $\mathcal G$ modulo translations applies.  Since
the limiting multiplier has degree zero, the value does not depend on
the kinetic normalization used for the equality manifold.

\begin{theorem}[Energy of a concentrating sequence]
\label{thm:r11-concentrating-sequence}
Suppose
\[
 M(f_n)\le M_c,\qquad A(f_n)\longrightarrow\infty,\qquad
 \sup_n E_k(f_n)<\infty.
\]
Then $M(f_n)\to M_c$.  Put $t_n=A(f_n)^{1/m}$ and write
$f_n(x)=t_n^{N/2}V_n(t_nx)$.  After extraction and translations,
$V_n\to Q$ strongly in $H^s$ for some $Q\in\mathcal G$, and
\begin{align}
 D(f_n)&\longrightarrow
   \int d_\infty(\xi/|\xi|)|\widehat Q(\xi)|^2\,d\xi,
       \label{eq:r11-deficit-limit}\\
 \liminf_n E_k(f_n)&\ge
   -\frac12\int d_\infty(\xi/|\xi|)|\widehat Q(\xi)|^2\,d\xi
   \ge e_\infty.
       \label{eq:r11-concentrating-liminf}
\end{align}
\end{theorem}

\begin{proof}
The sharp inequality and the mass bound give $E_a(f_n)\ge0$.
Boundedness of $d$ gives
\[
 0\le E_a(f_n)=E_k(f_n)+\tfrac12D(f_n)
      \le \sup_j E_k(f_j)+\tfrac12\|d\|_\infty M_c.
\]
Consequently
\[
 A(V_n)=1,\qquad
 E_a(V_n)=t_n^{-m}E_a(f_n)\longrightarrow0.
\]
Equation~\eqref{eq:r11-sharp-gn} also gives
\[
 0\le\frac12\left[1-\left(\frac{M(f_n)}{M_c}\right)^r\right]
      \le E_a(V_n)\longrightarrow0.
\]
Thus $M(f_n)\to M_c$, and normalized sharp-GN compactness yields
the asserted strong convergence.  Equivalently one may first correct
the masses and kinetic normalizations by factors tending to one and
then undo those normalizations after obtaining a compact subsequence.

Translations do not change Fourier modulus.  Since $d$ is bounded,
strong $L^2$ convergence implies
\[
 \left|\int d(t_n\xi)
      \bigl(|\widehat V_n(\xi)|^2-|\widehat Q(\xi)|^2\bigr)\,d\xi\right|
 \le\|d\|_\infty\|V_n-Q\|_2(\|V_n\|_2+\|Q\|_2)
 \longrightarrow0.
\]
For every $\xi\ne0$, \eqref{eq:r11-angular-limit} gives
$d(t_n\xi)\to d_\infty(\xi/|\xi|)$.  Dominated convergence now
proves~\eqref{eq:r11-deficit-limit}.  Finally
$E_k(f_n)=E_a(f_n)-D(f_n)/2\ge-D(f_n)/2$, proving
\eqref{eq:r11-concentrating-liminf}.
\end{proof}

\begin{corollary}[Bounded energy sublevels and sharpness]
\label{cor:r11-sublevels}
For every $e<e_\infty$,
\begin{equation}\label{eq:r11-uniform-sublevel}
 \sup\{A(f):M(f)\le M_c,\ E_k(f)\le e\}<\infty,
\end{equation}
where the assertion is vacuous if the set is empty.
Conversely, for every $e>e_\infty$ the exact-mass sublevel
$\{f:M(f)=M_c,\ E_k(f)\le e\}$ is unbounded in $H^s$.
\end{corollary}

\begin{proof}
Failure of~\eqref{eq:r11-uniform-sublevel} contradicts
Theorem~\ref{thm:r11-concentrating-sequence}.  For sharpness choose
$Q\in\mathcal G$ attaining the maximum in
\eqref{eq:r11-concentration-energy} and let
$Q_\lambda(x)=\lambda^{N/2}Q(\lambda x)$.  Homogeneity gives
\[
 M(Q_\lambda)=M_c,\qquad A(Q_\lambda)=\lambda^m,\qquad
 E_a(Q_\lambda)=0.
\]
Moreover
\[
 E_k(Q_\lambda)=-\tfrac12\int d(\lambda\xi)|\widehat Q(\xi)|^2\,d\xi
      \longrightarrow e_\infty
      \qquad(\lambda\longrightarrow\infty).
\]
For any fixed $e>e_\infty$ these trials eventually belong to the
stated sublevel while their kinetic energies diverge.
\end{proof}

The strict comparison in this result is optimal for mass--energy
coercivity.  At $e=e_\infty$, boundedness depends on the particular
symbol.  Sharpness of an energy sublevel does not itself construct a
blow-up trajectory.

\begin{remark}[Choice of the zero of frequency]
Replacing $k$ by $k+c$ multiplies solutions by $e^{-ict}$, adds $cM/2$
to their energy, and adds $cM_c/2$ to $e_\infty$.  Thus the difference
$E_k-e_\infty$ at $M=M_c$ is unchanged.  The concentration-energy
comparison does not depend on the arbitrary zero of frequency.
\end{remark}

\subsection{Conservation turns the bound into global dynamics}

For the general-symbol statement, we specify the local evolution input.
Let \textup{(F)} denote the following properties of
\begin{equation}\label{eq:r11-flow}
 i\partial_tu=k(D)u-|u|^{p-2}u:
\end{equation}
for every $u_0\in H^s$ there is a unique maximal solution
$u\in C((T_-,T_+),H^s)$, depending continuously on its initial datum,
conserving $M$ and $E_k$, and satisfying the continuation alternative
that a finite lifespan endpoint forces divergence of the $H^s$ norm.
Ellipticity of an arbitrary symbol is not by itself being used as a
substitute for this local evolution theory.

\begin{theorem}[Global bounded evolution at critical mass]
\label{thm:r11-global-endpoint}
Assume \textup{(F)}.  If
\[
 M(u_0)\le M_c,\qquad E_k(u_0)<e_\infty,
\]
then the solution of~\eqref{eq:r11-flow} is global in both time
directions and
\[
 \sup_{t\in\mathbb R}\|u(t)\|_{H^s}<\infty.
\]
For every strictly subcritical datum $M(u_0)<M_c$, the same conclusion
holds without any restriction on its energy.
\end{theorem}

\begin{proof}
In the first case, conservation places the maximal orbit in the fixed
sublevel of Corollary~\ref{cor:r11-sublevels} with
$e=E_k(u_0)<e_\infty$.  Ellipticity of $a$ and conservation of mass
bound its $H^s$ norm.  The continuation alternative excludes both
finite endpoints.  For the second assertion the direct estimate
\begin{equation}\label{eq:r11-subcritical-coercivity}
 E_k(f)\ge\frac12\left[1-\left(\frac{M(f)}{M_c}\right)^r\right]A(f)
             -\frac12\|d\|_\infty M(f)
\end{equation}
is coercive at each fixed mass below $M_c$; conservation gives the
same conclusion.
\end{proof}

\begin{proposition}[A concrete local theory including anisotropy]
\label{prop:r11-quartic-local}
For $m=4$ and $N\in\{1,2,3\}$, hypothesis \textup{(F)} holds for
every symbol $k=a-d$ in this section.  Consequently
Theorem~\ref{thm:r11-global-endpoint} is unconditional for these
quartic equations, including anisotropic $a$.
\end{proposition}

\begin{proof}
Here $s=2$ and the nonlinear map is $F(z)=|z|^\alpha z$ with
$\alpha=8/N>2$.  Its second derivative, viewing $\mathbb C$ as
$\mathbb R^2$, is locally Lipschitz.  The embeddings
$H^2(\mathbb R^N)\hookrightarrow L^\infty\cap W^{1,4}$ and the
ordinary chain and product rules imply, on every bounded $H^2$ ball,
\begin{equation}\label{eq:r11-local-lipschitz}
 \|F(u)-F(v)\|_{H^2}\le C_R\|u-v\|_{H^2}
 \quad\text{if }\|u\|_{H^2}+\|v\|_{H^2}\le R.
\end{equation}
For example the second-derivative terms have the forms
$F'(u)D^2u$ and $F''(u)(Du)^2$; their differences are bounded using
$H^2\hookrightarrow L^\infty$, $H^2\hookrightarrow W^{1,4}$,
and the local Lipschitz bounds on $F'$ and $F''$.

Since $k$ is real, $e^{-itk(D)}$ is a strongly continuous unitary group
on $H^2$.  Its Duhamel map and~\eqref{eq:r11-local-lipschitz} give a
contraction in $C([-T,T],H^2)$ with $T$ depending only on the initial
$H^2$ norm.  This proves existence, uniqueness, continuous dependence,
and the $H^2$ continuation alternative.

To justify conservation without assuming excessive smoothness of the
non-polynomial power, let $P_R$ be the orthogonal Fourier projection to
$\{\xi:|\xi|\le R\}$ and solve
\[
 i\partial_tu_R=k(D)u_R-P_RF(u_R),\qquad u_R(0)=P_Ru_0.
\]
The band-limited subspace is preserved.  On that subspace $k(D)$ is a
bounded operator on $H^2$, so the solution is differentiable in $H^2$.
The real $L^2$ pairings defining the derivatives of mass and energy are
therefore legitimate.  If $g_R=k(D)u_R-F(u_R)$, then
$\partial_tu_R=-iP_Rg_R$, so
\[
 \frac{d}{dt}E_k(u_R)
   =\operatorname{Re}\langle g_R,-iP_Rg_R\rangle=0.
\]
Likewise $dM(u_R)/dt=0$, since $P_Ru_R=u_R$ and
$\int\overline{u_R}F(u_R)$ is real.  The local contraction bounds are
uniform in $R$.  Strong convergence $P_R\to\mathrm{Id}$ on $H^2$
and the Duhamel equations yield
$u_R\to u$ in $C([-T,T],H^2)$.  Continuity of $M$ and $E_k$ passes
conservation to $u$.  Iterating over the maximal lifespan proves
\textup{(F)}.
\end{proof}

\subsection{Consequences for the two-threshold interpolation}

Let $a_0$ and $a_\infty$ be positive elliptic real-even symbols of degree
$m$, with $h=a_\infty-a_0$ positive elliptic.  Consider
\begin{equation}\label{eq:r11-rational-symbol}
 \begin{gathered}
 k_q(\xi)=a_0(\xi)+h(\xi)\frac{|\xi|^q}{1+|\xi|^q}
         =a_\infty(\xi)-d_q(\xi),\\
 d_q(\xi)=\frac{h(\xi)}{1+|\xi|^q},\qquad q>0.
 \end{gathered}
\end{equation}
Denote the two homogeneous sharp masses by $M_0<M_\infty$.
The following existence and compactness inputs are imported from the
companion study \emph{Thresholds, spectral geometry, and persistent
radiation in mass-critical dispersive equations}, Section~5; they are
not reproved in the present dynamical argument.  That analysis gives
$I(M)=\inf_{M(f)=M}E_{k_q}(f)<0$ with attainment for
$M_0<M<M_\infty$.  At $M=M_\infty$, the infimum is also negative
and attained when $q>m$, and every minimizing sequence is compact in
$H^s$ modulo translations.

\begin{corollary}[All negative-energy endpoint data are bounded]
\label{cor:r11-rational-global}
Assume a conservative local flow with the properties \textup{(F)} for
$k_q$.  For every $q>0$, all initial data of mass less than
$M_\infty$ generate global uniformly $H^s$-bounded solutions.
If $q>m$, the same conclusion holds for every initial datum satisfying
\begin{equation}\label{eq:r11-negative-endpoint-class}
 M(u_0)=M_\infty,\qquad E_{k_q}(u_0)<0.
\end{equation}
For $m=4$, $N=1,2,3$, and $q>4$, the latter assertion requires no
additional local-flow hypothesis.
\end{corollary}

\begin{proof}
For all $q>0$, and every $\eta>0$, homogeneity gives a finite
$C_\eta$ such that
$k_q\ge(1-\eta)a_\infty-C_\eta$.  Hence
\[
 E_{k_q}(f)\ge\frac12
 \left[1-\eta-\left(\frac{M(f)}{M_\infty}\right)^r\right]A_\infty(f)
       -\frac12C_\eta M(f).
\]
For a fixed subcritical mass choose $\eta$ to make the coefficient
positive and use conservation.  If $q>m$, $d_q$ is bounded and tends
uniformly to zero at high frequency.  Thus $d_\infty=0$ and
$e_\infty=0$, so Theorem~\ref{thm:r11-global-endpoint} proves
\eqref{eq:r11-negative-endpoint-class}.  The final assertion follows
from Proposition~\ref{prop:r11-quartic-local}.
\end{proof}

The endpoint result covers arbitrary data in the stated class: it uses
no radial symmetry, momentum restriction, or small distance to a
standing wave.  It excludes both finite-time blow-up and unbounded
infinite-time growth of the energy-space norm.

For $q=m$, the bounded correction has angular limit $h(\theta)$, so
\eqref{eq:r11-concentration-energy} supplies the precise endpoint
comparison.  If $I(M_\infty)<e_\infty$, its strict sublevel is nonempty
and Theorem~\ref{thm:r11-global-endpoint} applies there.  In the
isotropic case $h(\xi)=\varepsilon|\xi|^m$,
\[
 e_\infty=-\tfrac12\varepsilon M_\infty,
 \qquad
 E_{k_m}(f)>-\tfrac12\varepsilon M_\infty
 \quad\text{for every nonzero }f\text{ of mass }M_\infty.
\]
Thus that strict sublevel is empty; no endpoint global conclusion is
being inferred in the isotropic borderline case.  For $q<m$ the
correction is unbounded at high frequency and the concentration theorem
above does not apply.  The variational equality
$I(M_\infty)=-\infty$ in that case does not itself prove finite-time
blow-up.

For $q>m$, if $Q$ is an $a_\infty$ equality profile, then
$E_{k_q}(Q_\lambda)<0$ for every finite $\lambda$, while
$E_{k_q}(Q_\lambda)\to0$ and $A_\infty(Q_\lambda)\to\infty$.
A fixed negative energy therefore gives a bound, but no uniform bound
can cover all negative energies approaching zero.

\subsection{Stability of minimizing sets and the critical mass boundary}

At a mass $c$ where the constrained infimum is attained define
\[
 \Sigma_c=\{U\in H^s:M(U)=c,\ E_k(U)=I(c)\}.
\]
This entire set includes translations and constant phases.  Its
stability makes no uniqueness assertion about individual profiles.

\begin{theorem}[Stability in the correct perturbation class]
\label{thm:r11-set-stability}
Assume \textup{(F)}.
\begin{enumerate}
 \item For the interpolation~\eqref{eq:r11-rational-symbol}, if
 $M_0<c<M_\infty$, then $\Sigma_c$ is orbitally stable under all
 sufficiently small $H^s$ perturbations.
 \item For a bounded correction as above, suppose $c=M_c$,
 $I(c)<e_\infty$, and every minimizing sequence at mass $c$ is
 precompact in $H^s$ modulo translations.  Then $\Sigma_c$ is orbitally
 stable under sufficiently small $H^s$ perturbations whose masses are
 at most $c$.  In particular it is stable on the exact mass surface
 $M=c$.  This applies to~\eqref{eq:r11-rational-symbol} at
 $c=M_\infty$ whenever $q>m$.
\end{enumerate}
In both assertions all sufficiently close initial data in the indicated
perturbation class generate global solutions.  More explicitly, for
any $\epsilon>0$ there is $\delta>0$ such that
$\operatorname{dist}_{H^s}(u_0,\Sigma_c)<\delta$ implies
\[
 \sup_{t\in\mathbb R}\operatorname{dist}_{H^s}(u(t),\Sigma_c)<\epsilon,
\]
with the additional mass restriction in the second assertion.
\end{theorem}

\begin{proof}
The compactness assumption implies that $\Sigma_c$ is bounded in
$H^s$.  Thus closeness to this set makes the mass and energy uniformly
close to $c$ and $I(c)$, respectively.  In the first assertion all
sufficiently close data have mass uniformly below $M_\infty$; the
subcritical coercivity estimate used in
Corollary~\ref{cor:r11-rational-global} makes their solutions global
and bounds their $H^s$ norms uniformly.  In the second assertion choose
$e$ strictly between $I(c)$ and $e_\infty$ and apply
Corollary~\ref{cor:r11-sublevels} to masses at most $c$ and energies
at most $e$.  This gives the same global existence and uniform bound.

If stability failed, there would be initial data $u_{0,n}$ with
$\operatorname{dist}_{H^s}(u_{0,n},\Sigma_c)\to0$ and times $t_n$
such that $f_n=u_n(t_n)$ stays a fixed positive distance from
$\Sigma_c$.  Conservation gives
\[
 M(f_n)\longrightarrow c,\qquad E_k(f_n)\longrightarrow I(c),
\]
and the preceding estimates bound $f_n$ in $H^s$.  Set
$g_n=(c/M(f_n))^{1/2}f_n$.  Then
\[
 M(g_n)=c,\qquad \|g_n-f_n\|_{H^s}\longrightarrow0,\qquad
 E_k(g_n)\longrightarrow I(c).
\]
Compactness of exact-mass minimizing sequences now gives
\[
 \operatorname{dist}_{H^s}(g_n,\Sigma_c)\longrightarrow0.
\]
The same follows for $f_n$, a contradiction.
\end{proof}

The proof is the established variational compactness mechanism for
orbital stability.  The additional concentration-energy bound is what
allows it to operate at the critical mass.  At that boundary the
restriction $M(u_0)\le M_c$ cannot be silently omitted: nearby data can
otherwise cross into the supercritical mass range, where the argument
supplies no uniform kinetic bound.  Neither unconstrained endpoint
stability nor instability is asserted here.

\subsection{Removal of the flat scale direction}

The interpolation also produces an exact measure of the energy cost of
changing the spatial scale of a minimizer.  This statement does not
require spectral nondegeneracy.

\begin{proposition}[Exact scale curvature]
\label{prop:r11-scale-curvature}
Let $U\ne0$ be a constrained energy minimizer for
\eqref{eq:r11-rational-symbol}.  Define
\[
 U_t(x)=e^{Nt/2}U(e^tx),\qquad
 \mathcal K(U)=\left.\frac{d^2}{dt^2}E_{k_q}(U_t)\right|_{t=0}.
\]
Then the scalar energy curve is twice differentiable and, with
$z=|\xi|^q$,
\begin{equation}\label{eq:r11-exact-curvature}
 \mathcal K(U)=\frac q2\int
 h(\xi)\frac{z[(m+q)+(m-q)z]}{(1+z)^3}
                    |\widehat U(\xi)|^2\,d\xi.
\end{equation}
For $0<q\le m$, $\mathcal K(U)>0$.  For $q>m$, minimality still
gives $\mathcal K(U)\ge0$, but positivity does not follow pointwise
from the integrand.
\end{proposition}

\begin{proof}
Let $b(\xi)=h(\xi)|\xi|^q/(1+|\xi|^q)$ and let
$\mathcal D=\xi\cdot\nabla_\xi$.  The derivatives
$\mathcal D^jb$, $j=0,1,2$, are bounded in absolute value by
$C|\xi|^m$.  Therefore differentiating the scalar Fourier energy is
justified by dominated convergence, with no assumption of a spatial
moment or differentiability of $t\mapsto U_t$ in $H^s$.  Homogeneity
gives
\[
 E_{k_q}(U_t)=e^{mt}E_{a_0}(U)
                 +\frac12\int b(e^t\xi)|\widehat U(\xi)|^2\,d\xi.
\]
Its first derivative vanishes at the scalar minimum $t=0$.  Eliminating
$E_{a_0}(U)$ from the second derivative yields
\[
 \mathcal K(U)=\frac12\int(\mathcal D^2-m\mathcal D)b(\xi)
                                    |\widehat U(\xi)|^2\,d\xi.
\]
Since $\mathcal Dh=mh$ and $\mathcal Dz=qz$, direct differentiation
gives
\[
 (\mathcal D^2-m\mathcal D)b
      =q h\frac{z[(m+q)+(m-q)z]}{(1+z)^3},
\]
proving~\eqref{eq:r11-exact-curvature}.  If $q\le m$, its coefficient
is strictly positive for every $\xi\ne0$, and a nonzero $L^2$ function
cannot have all Fourier mass at the origin.  If $q>m$, the coefficient
changes sign at $z=(m+q)/(q-m)$; nonnegativity of the integral follows
instead from minimality.
\end{proof}

For the homogeneous critical optimizer the energy is exactly flat
under mass-preserving dilation.  Proposition~\ref{prop:r11-scale-curvature}
proves that the rational correction removes this degeneracy in the
scale direction when $q\le m$.  It does not identify a complete
Hamiltonian spectral mode.

We give the precise lower-threshold consequence of the selection theorem
established for this interpolation.  Put
\[
 \begin{gathered}
 \alpha=\frac{M}{M_0},\qquad d=\alpha^r-1,\qquad
 U_M(x)=\sqrt\alpha\,\lambda_M^{N/2}V_M(\lambda_Mx),\\
 M(V_M)=M_0,\qquad A_0(V_M)=1.
 \end{gathered}
\]
Let
\[
 \begin{gathered}
 \mathcal G_0=\{Q:M(Q)=M_0,\ A_0(Q)=1,\ E_{a_0}(Q)=0\},\\
 C(Q)=\int h(\xi)|\xi|^q|\widehat Q(\xi)|^2\,d\xi,
 \end{gathered}
\]
and suppose at least one element of $\mathcal G_0$ has finite $C$.
Write $C_{\min}=\min_{\mathcal G_0}C$ and
\[
 \beta_*^q=\frac{m}{(m+q)C_{\min}},\qquad
 K_* =\frac{q}{2(m+q)}\beta_*^m.
\]
The lower-threshold selection theorem of the same companion study
provides the following inputs; the new deduction here is the curvature
asymptotic, not a new proof of profile selection:
\begin{align}
 \lambda_M&\sim\beta_*d^{1/q},&
 I(M)&\sim-\alpha K_*d^{1+m/q},
       \label{eq:r11-selection-recall-a}\\
 \frac{B(U_M)}{A_0(U_M)}&=\frac{m}{m+q}d+o(d),&
 \int\frac{h(\xi)|\xi|^q}{1+|\lambda_M\xi|^q}
                |\widehat V_M(\xi)|^2\,d\xi&\longrightarrow C_{\min},
       \label{eq:r11-selection-recall-b}
\end{align}
where $B(f)=\int b|\widehat f|^2$.  After translations and extraction,
$V_M$ converges strongly in $H^s$ to some $Q\in\mathcal G_0$ with
$C(Q)=C_{\min}$.  The correction-energy convergence in
\eqref{eq:r11-selection-recall-b}, not merely strong $H^s$
convergence, is needed in the next deduction.

\begin{corollary}[Sharp lower-threshold scale curvature]
\label{cor:r11-curvature-asymptotic}
For minimizers satisfying the preceding selection conclusions, for
every $q>0$,
\begin{equation}\label{eq:r11-curvature-asymptotic}
 \mathcal K(U_M)\sim\frac{mq}{2}dA_0(U_M),\qquad
 \frac{\mathcal K(U_M)}{|I(M)|}\longrightarrow m(m+q)
 \qquad(M\downarrow M_0).
\end{equation}
In particular the scale curvature is strictly positive near the lower
threshold also when $q>m$.
\end{corollary}

\begin{proof}
Along a convergent selected-profile subsequence define the nonnegative
finite measures
\[
 d\mu_M(\xi)=
 \frac{h(\xi)|\xi|^q}{1+|\lambda_M\xi|^q}
                       |\widehat V_M(\xi)|^2\,d\xi.
\]
Their total masses tend to $C_{\min}$.  On every compact annulus they
converge to
$h(\xi)|\xi|^q|\widehat Q(\xi)|^2\,d\xi$, whose total mass is also
$C_{\min}$.  Positivity therefore gives tightness away from both zero
and infinity: exhaust the limiting integral by compact annuli and use
total-mass convergence to control the complements.

Formula~\eqref{eq:r11-exact-curvature} becomes
\[
 \frac{\mathcal K(U_M)}{\alpha\lambda_M^{m+q}}
 =\frac q2\int
 \frac{(m+q)+(m-q)|\lambda_M\xi|^q}
      {(1+|\lambda_M\xi|^q)^2}\,d\mu_M(\xi).
\]
The displayed multiplier is uniformly bounded and tends on compact
annuli to $m+q$.  Tightness proves that the right side converges to
$q(m+q)C_{\min}/2$.  Since
$B(U_M)=\alpha\lambda_M^{m+q}\mu_M(\mathbb R^N)$, we have
\[
 \frac{\mathcal K(U_M)}{B(U_M)}\longrightarrow\frac{q(m+q)}2.
\]
Combining this with~\eqref{eq:r11-selection-recall-b} proves the first
assertion.  The energy and scale laws in
\eqref{eq:r11-selection-recall-a}, and
$A_0(U_M)=\alpha\lambda_M^m$, give the second.  Every selected
subsequence has the same scalar limits, so the conclusions hold along
the full family.
\end{proof}

There is also a fixed-scale version.  Let
$b_0(\xi)=a(\xi)|\xi|^q/(1+|\xi|^q)$,
$k_\varepsilon=a+\varepsilon b_0$, and fix $0<\tau<1$.
For the joint limit with
$\alpha_\varepsilon=(1+\varepsilon\tau)^{1/r}$, suppose the
selected minimizers satisfy
$U_\varepsilon/\sqrt{\alpha_\varepsilon}\to Q$ strongly in $H^s$
modulo translations.  Since
$(\mathcal D^2-m\mathcal D)b_0$ is an order-$m$ multiplier,
\begin{equation}\label{eq:r11-joint-curvature}
 \frac{\mathcal K(U_\varepsilon)}
      {\alpha_\varepsilon\varepsilon}
 \longrightarrow\frac12\int
       (\mathcal D^2-m\mathcal D)b_0(\xi)|\widehat Q(\xi)|^2\,d\xi.
\end{equation}
For $q\le m$ the limit is positive.  Thus an arbitrarily weak
perturbation creates a scale curvature of order $\varepsilon$ at the
finite spatial scale selected by a mass excess of the same order.

An oscillation frequency requires a further Hamiltonian spectral
analysis.  The conjugate phase, translation coupling, and interaction
with other spectral modes enter that calculation.  The exact curvature
and its asymptotics above do not presume that such an isolated
oscillatory eigenmode has already been established.

\subsection{Which parts extend to other critical interactions}

The concentration argument uses only a homogeneous energy
$E_a=A/2-\mathcal F$ with $E_a\ge0$ at mass at most $M_c$,
compactness of normalized zero-energy sequences, and a bounded
quadratic correction $-D/2$.  It does not use locality or a polynomial
formula for $\mathcal F$.  Theorems~\ref{thm:r11-concentrating-sequence}
and~\ref{thm:r11-global-endpoint} consequently extend, with the same
proof, to a mass-critical Hartree interaction once its sharp inequality,
optimizer compactness, and conservative local flow have been verified.
The curvature calculation needs only
$\mathcal F(U_t)=e^{mt}\mathcal F(U)$, in addition to the explicit
dispersion.  Set stability further requires compactness of the relevant
constrained minimizing sequences.  This separates the common
concentration mechanism from the interaction-specific analytic inputs.
